\section*{Part VI: Hiring Managers and Executives}

\subsection*{62. A Manager's Principal Product Is Other People's Performance}

Promotion often rewards the strongest individual contributor, but management requires a different production function: allocating work, setting standards, coaching, giving difficult feedback, coordinating dependencies, resolving conflict, and developing successors. Being technically brilliant does not automatically predict being a strong manager. Nor does a warm interpersonal style establish willingness to make hard trade-offs.

Evaluate managers on \textbf{team diagnosis, delegation, performance management, operational systems, judgment under conflict, and the creation of capability in others}. Require evidence from employees or colleagues, not only the manager's own account. A manager who says ``my team delivered'' should be able to explain how goals were set, what they delegated, which decisions employees made independently, how underperformance was addressed, and what improved without their direct involvement.

\subsection*{63. Management Simulation: The Talented but Unreliable Contributor}

Give a candidate a case: a high-performing engineer produces exceptional output but repeatedly fails to document decisions, interrupts colleagues, and misses cross-team commitments. Other contributors report increasing frustration, while a launch deadline approaches. Ask for the next conversation, a 30-day plan, success measures, and an escalation path.

A weak manager avoids the conflict, threatens immediate dismissal without investigation, or says the star should be exempt from standards. A strong 1 gathers evidence, distinguishes capability from behavior, specifies expectations, offers support and a timeline, addresses team impact, and follows applicable policy. They neither protect performance at any cost nor enforce rules without judgment. For a learning institution, the equivalent could be a celebrated instructor who ignores assessment rubrics or belittles struggling students.

\subsection*{64. Evaluate Delegation Using Real Constraints}

Ask the candidate to lead a fictional team of 3 with a 4-week deliverable: 1 expert is part time, 1 junior needs mentorship, a key dependency is outside the team, and the founder frequently changes priorities. The candidate must write a plan that identifies outcomes, decision owners, cadence, checkpoints, and what is explicitly not being done.

Strong managers protect focus while leaving individuals room for judgment. They define escalation thresholds, avoid micromanagement, and explain the costs of coordination. Candidates who personally volunteer to complete every task may be energetic individual contributors but have not demonstrated managerial leverage. For a senior manager, require a plausible plan to make the system function during their 2-week absence.

\subsection*{65. Managerial Evidence Should Include Downward and Lateral Views}

A supervisor reference can confirm targets, but direct reports and peers often observe different aspects of leadership: psychological safety, credit allocation, quality of feedback, promises kept, and treatment of disagreement. With permission and appropriate privacy safeguards, seek structured examples from more than 1 vantage point. Do not solicit unauthorized confidential data or treat an unverified accusation as a fact; give the candidate an opportunity to clarify material inconsistencies.

Ask ``What do people on this person's team become better at?'' and ``What problems repeatedly required escalation?'' Such questions move beyond popularity. Some effective managers make necessary unpopular decisions; some beloved managers avoid performance accountability. References should be interpreted in the context of the role and organizational culture, not reduced to a simple ``liked by team'' count.

\subsection*{66. The Manager Scorecard}

An illustrative functional-manager rubric is:
\begin{itemize}[leftmargin=*]
\item Team development and coaching: 25\%
\item Execution systems and process design: 20\%
\item Judgment and prioritization: 20\%
\item Feedback, accountability, and conflict resolution: 15\%
\item Cross-functional collaboration: 10\%
\item Technical or domain credibility: 10\%
\end{itemize}
with integrity, respect, and essential legal duties as non-negotiable gates. Compare this with an individual-contributor scorecard: technical expertise often remains necessary but no longer dominates every decision.

For a 1st-line technical lead who still codes, adjust the mix. For a program director coordinating instructors across colleges, emphasize instructional quality assurance and contract delivery. For a general manager, include financial ownership and resource allocation. A founder should resist purchasing a prestigious management title without defining which decisions the manager will be authorized to make.

\subsection*{67. Executive Hiring Is a Theory of Decisions Under Ambiguity}

Executives work on a larger time scale and make decisions with consequences that often cannot be quickly reversed: strategy, capital, organizational structure, leadership succession, external trust, and institutional commitments. The best evidence is not a polished future vision in isolation but a \textbf{track record of difficult choices and how the executive reasoned, learned, and responded after outcomes emerged}.

Examine resource allocation, strategic clarity, ability to build high-quality teams, governance, ethical conduct, response to failure, and capacity to adapt to an unfamiliar context. An executive's prior success may owe heavily to a brand, abundant funding, exceptional peers, market timing, or inherited scale. The founder must distinguish those enabling conditions from portable judgment. It is especially dangerous to hire an executive simply to confer legitimacy on a newborn enterprise.

\subsection*{68. The Executive Capital-Allocation Exercise}

Give a hypothetical CEO, CFO, or COO finalist a scenario: \text{₹}1~crore of available discretionary capital over 2 years; potential investments in customer acquisition, instructor capacity, platform automation, a research pilot, and a modest physical facility; current revenue uncertain; a liquidity reserve requirement; and different stakeholder claims. Require an allocation plan with explicit assumptions, downside cases, opportunity costs, decision gates, and what to stop funding.

A strong executive connects timing, risk, reversibility, mission, and measurable outcomes. They reserve capital for obligations, distinguish venture experiments from repeatable operations, and recognize that attractive long-term initiatives can be premature. A weak executive allocates everything to the most visible expansion project or uses optimistic forecasts without stress tests. The test should not reward agreement with the founder's preferred investment: \textbf{thoughtful dissent is evidence}.

\subsection*{69. The Executive Personnel Decision Case}

Provide a scenario in which a trusted senior leader is producing growth while concealing quality failures or pressuring staff to manipulate success metrics. Ask how the candidate would investigate, protect due process, preserve evidence, communicate with stakeholders, and decide on corrective action. This tests ethical leadership, governance, and willingness to confront inconvenient facts.

It is inappropriate to infer moral character from 1 scripted answer. The case is a prompt to seek corroborating past behavior and clear commitments. Conduct lawful references and due diligence, particularly for senior roles entrusted with funds, students, confidential information, or public claims. Do not mistake an assertive tone for courage, or a diplomatic tone for lack of backbone.

\subsection*{70. The Executive's Attraction to Prestige Can Be a Warning}

A senior candidate may be interested in the promised future institution, international affiliations, campus, or research status while uninterested in the current work of building small, repeatable, unglamorous systems. That mismatch is decisive for an early venture. A founder should describe the actual state of operations, financial limits, authorization, reporting lines, and expected hands-on contribution before inviting an executive to commit.

The candidate should also evaluate whether the job is responsible: can the founder delegate real authority, support the role, and accept disagreement? An experienced operator may rightly decline a title with responsibility but no control. Great hiring is mutual selection. An inflated job description is not merely marketing excess; it is a forecast of future conflict.

\subsection*{71. Executive Interview: 6 Decisions from the Candidate's Past}

Ask for:
\begin{enumerate}[leftmargin=*]
\item 1 major investment approved,
\item 1 major investment rejected,
\item 1 costly mistake and recovery,
\item 1 team rebuilt or reorganized,
\item 1 strategic premise changed based on evidence, and
\item 1 ethical conflict handled.
\end{enumerate}
For each request the original information available, options considered, stakeholder influence, decision rule, evidence discovered later, and the candidate's actual contribution. Where appropriate, ask references about the same episodes and compare accounts respectfully.

Be alert to hindsight narratives in which every success was foresight and every failure was someone else's fault. The better executive often has a coherent explanation of what they did not know and which decision process they changed afterward. Successful outcomes matter; so does evidence that the process can survive a different environment. The best executive is not a person who has never been wrong, but 1 whose organization becomes better at discovering and correcting error.

\subsection*{72. The Executive Scorecard and Evidence Horizon}

A proposed senior leader rubric is:
\begin{itemize}[leftmargin=*]
\item Strategic judgment and capital allocation: 25\%
\item Organizational architecture and people development: 20\%
\item Execution under constraints: 20\%
\item Governance and ethical integrity: 15\%
\item External stakeholder effectiveness: 10\%
\item Learning and adaptation: 10\%
\end{itemize}
For a CFO, elevate cash discipline, financial controls, and reporting; for a chief academic officer, academic standards, faculty development, and learner-outcome integrity; for a CRO, repeatable ethical revenue and sales operating systems.

Assess across at least 2 horizons: what the executive can improve in 90 days and what institution-building changes they can credibly initiate over 2 years. Avoid offering equity or broad authority simply because the résumé is impressive. The higher the role's impact and irreversibility, the more important a documented decision memo, reference triangulation, governance clarity, and realistic downside analysis become.

\subsection*{73. Founder and Executive Fit Must Include Productive Disagreement}

The founder's mission can be a strategic boundary without becoming a demand for ideological conformity. A senior hire should share non-negotiable standards and institutional purpose while challenging assumptions about market, timing, capital, pedagogy, and hiring. A candidate who expertly mirrors the founder's vocabulary may be good at interviews rather than governance.

Test a live disagreement: ``Here is our proposed expansion plan. What would you stop? Which assumption do you think is most dangerous? What would convince you that you were wrong?'' Reward accurate reasoning, evidence, respect, and ability to commit after a decision. Neither habitual contrarianism nor automatic assent is desirable. What matters is making the organization's decisions better.
